\begin{corollary} \label{17:regular_to_int}
    Функция регулярная в односвязной области имеет первообразную в этой области.
\end{corollary}

\begin{proof}
    Пусть \(f\) регулярна в односвязной области \(G\) и возьмём КНД контур \(\gamma \subset G\).
    Тогда \(\int\limits_\gamma f dz = 0\). Так как регулярная функция непрерывна, то из утверждения
    \ref{12:int_kontur} вытекает, что \(f\) имеет первообразную.
\end{proof}

\begin{note}
    Следствие \ref{17:regular_to_int} неверно для неодносвязной области.
\end{note}

\begin{example}
    \(f(z) = e^{\ctg{z}}\) имеет первообразную в верхней полуплоскости комплексного пространства
     \(\{z: \im z > 0\}\).
\end{example}

\begin{definition}
    Область \(G \subset \R^2\) называется областью с простой границей (ОСПГ), если \(\partial G =
    \gamma_1 \dots \gamma_n\), где \(\gamma_1, \dots, \gamma_n\)~--- простые попарно не 
    пересекающиеся контуры. \(\partial G\) будем называть положительно ориентированной, если при 
    обходе любой \(\gamma_k, \, k \in \overline{1,n}\) область \(G\) остаётся слева.
\end{definition}

\begin{example}
    Множество \(B_1(0)\) с границей в виде окружности является ОСПГ. Множество \(\mathring{B}_1(0)\)
    не является ОСПГ, так как её граница представима в виде окружности и точки (точку контуром не 
    считаем). Множество \(\{z: \, \im z > 0\}\) не является ОСПГ, так как её граница~--- прямая.
    Множество \(G = \{z \in \Co \colon |z| > 1\} \cup \{\infty\}\) является ОСПГ. При этом множество
    \(G = \{z \in \Co \colon |z| > 1\} \) не является ОСПГ.
\end{example}

\begin{theorem} \label{17:ma_2}
    (МА). Пусть \(D, G \subset \R^2\)~--- области, причём \(G\)~--- ограниченная ОСПГ и 
    \(\overline{G} \subset D\), а также \(P, Q \in C^1(D)\). Тогда выполняется следующее равенство:
    \[\int\limits_{\partial G} Pdx + Qdy = \iint\limits_G (\dfrac{\partial Q}{\partial x} - 
    \dfrac{\partial P}{\partial y})dxdy\]
\end{theorem}

\begin{statement}\label{17:int_koshi_ogr_ospg}
    \textbf{(Интегральная теорема Коши для ограниченной ОСПГ).} Пусть \(G \subset \Co\)~--- 
    ограниченная ОСПГ и \(f\) регулярна в \(\overline{G}\). Тогда \(\int\limits_{\partial G} f dz = 
    0\).
\end{statement}

\begin{proof}
    Пусть \(f(z) = u(x, y) + i v(x, y)\). Тогда \(\int\limits_{\partial G} f dz = 
    \int\limits_{\partial G} (udx + vdy) + i \int\limits_{\partial G} (vdx + udy) = I_1 + iI_2\).
    По теореме  \ref{17:ma_2} для области \(G\), \(P=u\), \(Q = -v\) выполняется \(I_1 = 
    \int\limits_{\partial G} (udx + vdy) = \iint\limits_G (-v_x -u_y)dxdy \equiv 0\) (здесь учтено
    \(u_y = -v_x\) из утверждения \ref{9:Koshi_Riman}). Аналогично с \(I_2\). Таким образом 
    \(\int\limits_{\partial G} f dz = 0 + i \cdot 0 = 0\).
\end{proof}

\begin{example}
    Пусть \(f\) регулярна в замыкании кольца \(r_1 < |z-z_0| < r_2\). Тогда \(\int\limits_{C_1} 
    \dfrac{f(z) dz}{(z-z_0)^n} = \int\limits_{C_2} \dfrac{f(z) dz}{(z-z_0)^n}\), здесь \(C_1, C_2\)
    ~--- границы кольца с радиусами \(r_1\) и \(r_2\) соответственно с одинаковыми ориентациями.
\end{example}

\begin{proof}
    Обозначим подынтегральную функцию \(g(z) = \dfrac{f(z)}{(z-z_0)^n}\). Она является регулярной в 
    замыкании этого кольца (\(f\)~--- регулярна в замыкании кольца, а выражение в знаменателе не 
    обращается в нуль на множестве интегрирования), при этом само кольцо~--- ограниченная ОСПГ. 
    Тогда из утверждения \ref{17:int_koshi_ogr_ospg} следует, что \(\int\limits_{C_2 \cup (C_1)^{-1}
    } g(z) dz = 0\). Раскрыв данное выражение по свойствам интеграла вдоль кривой второго рода 
    получим доказываемое равенство.
\end{proof}


\Section{Интегральная формула Коши для ограниченной ОСПГ. Формулы для производных}

\begin{statement}\label{18:int_formula}
    Пусть \(G\)~--- ограниченная ОСПГ, \(f\) регулярна в \(\overline{G}\) и \(z \in G\). Тогда верны
    формулы:
    \[f^{(n)}(z) = \dfrac{n!}{2 \pi i} \int\limits_{\partial G} 
    \dfrac{f(\zeta) d\zeta}{(\zeta - z)^{n+1}}, \, n \in \overline{0, 1, \dots}.\]
\end{statement}

\begin{proof}
    Зафиксируем \(z \in G\). Возьмём любой прямоугольник \(\Pi\), для которого выполняется \(z \in 
    \Pi \subset \overline{\Pi} \subset G\). Тогда из следствия \ref{16:inf_dif_reg_func} следует
    \(f^{(n)}(z) = \dfrac{n!}{2 \pi i} \int\limits_{\partial \Pi} \dfrac{f(\zeta) d\zeta}
    {(\zeta - z)^{n+1}}\). Заметим, что \(G_1 = G \setminus \overline{\Pi}\) тоже ограниченная ОСПГ.
    Для неё, вследствие регулярности в \(G_1\) подынтегральной функции \(g(z)\) доказываемой 
    формулы, выполняется утверждение \ref{17:int_koshi_ogr_ospg}, то есть \(0 = 
    \int\limits_{\partial G_1} g(z) dz = \int\limits_{\partial G} g(z) dz + 
    \int\limits_{\partial \Pi ^{-1}} g(z) dz = \int\limits_{\partial G} g(z) dz - f^{(n)}(z)\). 
    Переносом слагаемых получаем доказываемую формулу.
\end{proof}

\begin{corollary}
    Если \(f(z)\) регулярна в круге \(K = {z \in \Co \colon |z-z_0| < R}\), то 
    \(f(z) = \sum\limits_{n = 0}^\infty \dfrac{f^{(k)}(z_0)}{k!} (z-z_0)^n\).
\end{corollary}

\begin{proof}
    Зафиксируем \(R_1 \in (0, R)\). Так как \(f(z)\) регулярна в круге \(K'\) \(|z-z_0| \le R_1\), 
    то по утверждению \ref{18:int_formula} \(\forall z \in K' \rightarrow f(z) = \dfrac{1}{2 \pi i} 
    \int\limits_{|z-z_0|=R_1} \dfrac{f(\zeta) d\zeta}{(\zeta - z)}\). Тогда из 
    утверждения \ref{14:int_koshi} \(f(z)\) представима степенным рядом.
\end{proof}


\Section{Гармонические функции}

\begin{definition}
    Пусть \(G \subset \R^2\)~--- область. \(u(x, y) \colon G\rightarrow\R\) называется гармонической
    в \(G\), если \(u(x, y) \in C^2(G)\) и \(u_{xx} + u_{yy} = \Delta u \equiv 0\).
\end{definition}

\begin{statement}
    Пусть \(f(z) = u(x, y) + iv(x, y)\) регулярна в \(G\). Тогда \(u, v\) гармонические в \(G\).
\end{statement}

\begin{proof}
    Пусть \(f\) регулярна в \(G\), тогда \(f \in C^\infty(G)\) (по следствию 
    \ref{16:inf_dif_reg_func}). Тогда \(u, v \in C^\infty(G)\). Из утверждения \ref{9:Koshi_Riman}
    имеем \(u_x = v_y\) и \(u_y = -v_x\), а значит \(u_{xx} = v_{yx}\) и \(u_{yy} = -v_{xy}\).
    Из равенства \(v_{xy} = v_{yx}\) получаем, что \(u_{xx} + u_{yy} = 0\), что совпадает с 
    определением гармоничности \(u\). Аналогично доказывается для \(v\).
\end{proof}

\begin{statement}
    Пусть \(u(x, y)\)~--- гармоническая функция в односвязной области \(G \subset \R^2\). Тогда 
    существует \(f(z) = u(x, y) + iv(x, y)\) регулярная в \(G\). 
\end{statement}

\begin{proof}
    Для доказательства воспользуемся дополнением к теореме \ref{17:ma}: для \(\omega(x, y) = 
    \int\limits_{(x_0, y_0)}^{(x, y)} Pdx + Qdy\) справедливо \(P = \dfrac{\partial \omega}{\partial
    x}\) и \(Q = \dfrac{\partial \omega}{\partial y}\). Тогда \(v(x, y) = 
    \int\limits_{(x_0, y_0)}^{(x, y)} -u_ydx + u_xdy\). Восстановить \(v\) удаётся с точностью до 
    константы
\end{proof}

\begin{note}
    Односвязность области существенна. \(u = \ln{(x^2 + y^2)}\) гармоническая в \(\R^2 \setminus 
    \{(0, 0)\}\), однако \(v\) восстановить нельзя. Если бы \(v\) существовала, то \(f_z = u_x - 
    iu_y = \dfrac{2}{z}\). Но тогда \(f\) первообразная \(\dfrac{2}{z}\), получаем противоречие.
\end{note}